\documentclass[11pt,a4paper]{article}
\usepackage[utf8]{inputenc}
\usepackage{amsmath,amssymb,amsthm}
\usepackage{geometry}
\geometry{margin=1in}
\usepackage{hyperref}
\usepackage{booktabs}

\title{Smale's Mean Value Conjecture \\ and the Pandrosion Vanishing Identity}
\author{Ivan Besevic}
\date{April 2026}

\newtheorem{theorem}{Theorem}[section]
\newtheorem{lemma}[theorem]{Lemma}
\newtheorem{proposition}[theorem]{Proposition}
\newtheorem{corollary}[theorem]{Corollary}
\newtheorem{conjecture}[theorem]{Conjecture}
\newtheorem{remark}[theorem]{Remark}

\begin{document}
\maketitle

\begin{abstract}
We reformulate Smale's mean value conjecture as a bound on the Pandrosion quotient: $\min_{P'(c) = 0} |Q(0, c)|/|Q(0, 0)| \le (d - 1)/d$, where $Q(z, w) = P(w)/w$. The classical Lagrange--Sylvester identity $\sum 1/P'(\alpha_k) = 0$ (written here in Pandrosion notation as $\zeta_P(1) = 0$), specialised to a root at the origin, yields the constraint $1 = -\sum 1/(\alpha_k R'(\alpha_k))$, from which the triangle inequality gives $\min_k |\alpha_k R'(\alpha_k)| \le d - 1$. This is a root-side bound, not a direct critical-point bound. We identify the extremal polynomial $P = z^d - z$, which attains $\min_c |P(c)/c|/|P'(0)| = (d - 1)/d$ exactly for every $d$, and report a numerical sweep over $40{,}000$ random polynomials ($d = 2, \dots, 9$) with no violations. Section~5 describes what an eventual proof would still require; this paper supplies structural constraints, not a full proof.
\end{abstract}

\paragraph{Scope statement.} We do \emph{not} claim a proof of Smale's MVC for any $d \ge 4$. Section~2 derives a \emph{root-side} bound $\min_k |\alpha_k R'(\alpha_k)| \le d - 1$ from the classical Lagrange--Sylvester identity; Section~3 gives the critical-point identity $S(c) = |c R'(c)|$; Section~5 explains why these two ingredients alone do not close the MVC and what a completed argument would still need.

\section{Smale's conjecture in Pandrosion form}

\begin{conjecture}[Smale, 1981]
For any polynomial $P$ of degree $d \ge 2$ with $P(0) = 0$ and $P'(0) \ne 0$, there exists a critical point $c$ (zero of $P'$) such that:
\begin{equation}
\frac{|P(c)|}{|c \cdot P'(0)|} \;\le\; \frac{d - 1}{d}.
\end{equation}
\end{conjecture}

Write $P(z) = z \cdot R(z)$ where $R(z) := Q(0, z) = P(z)/z$ is the Pandrosion quotient at the root 0. The Smale ratio becomes:
\begin{equation}
\boxed{S(c) \;:=\; \frac{|R(c)|}{|R(0)|} \;=\; \frac{|Q(0, c)|}{|Q(0, 0)|} \;=\; \frac{\prod_{k \ge 2} |c - \alpha_k|}{\prod_{k \ge 2} |\alpha_k|}.}
\end{equation}

The conjecture states: $\min_{P'(c) = 0} S(c) \le (d - 1)/d$.

\section{The vanishing identity at the origin}

\begin{theorem}[Specialised vanishing identity]
For $P$ with $P(0) = 0$ and $P'(0) = 1$ (normalised), and roots $0 = \alpha_1, \alpha_2, \dots, \alpha_d$:
\begin{equation}
\boxed{1 \;=\; -\sum_{k=2}^{d} \frac{1}{\alpha_k\, R'(\alpha_k)},}
\end{equation}
where $R(z) = \prod_{k \ge 2} (z - \alpha_k)$ and $R'(\alpha_k) = \prod_{j \ne k, j \ge 2} (\alpha_k - \alpha_j)$.
\end{theorem}

\begin{proof}
The identity $\sum_{k=1}^d 1/P'(\alpha_k) = 0$ is the classical \emph{Lagrange--Sylvester partial fraction identity} for $1/P(z)$ when $\deg P \ge 2$ (cf.\ \cite[Ch.~5]{markushevich}); it follows from comparing the partial-fraction expansion $1/P(z) = \sum_k 1/(P'(\alpha_k)(z - \alpha_k))$ with the Laurent expansion $1/P(z) = O(1/z^d) = O(1/z^2)$ at infinity, the $1/z$ coefficient on the right being $-\sum_k 1/P'(\alpha_k)$. Specialised to $P(z) = z R(z)$ with $P'(0) = R(0) = 1$ and $P'(\alpha_k) = \alpha_k R'(\alpha_k)$ for $k \ge 2$, we get $1 + \sum_{k \ge 2} 1/(\alpha_k R'(\alpha_k)) = 0$. The Pandrosion-zeta presentation~$\zeta_P(1) = \sum 1/P'(\alpha_k) = 0$ is purely a notation; the underlying identity is classical. \qedhere
\end{proof}

\begin{corollary}[Root separation bound]
By the triangle inequality on the identity:
\begin{equation}
1 \;\le\; \sum_{k=2}^d \frac{1}{|\alpha_k|\, |R'(\alpha_k)|} \;\le\; (d - 1) \cdot \max_k \frac{1}{|\alpha_k|\, |R'(\alpha_k)|},
\end{equation}
hence:
\begin{equation}
\boxed{\min_{k \ge 2} |\alpha_k| \cdot |R'(\alpha_k)| \;\le\; d - 1.}
\end{equation}
\end{corollary}

This says: at least one non-zero root $\alpha_k$ satisfies $|\alpha_k| \cdot |R'(\alpha_k)| \le d - 1$. Since $R'(\alpha_k)$ measures the separation of $\alpha_k$ from the other non-zero roots, this constrains how ``spread out'' the roots can be.

\section{The critical point identity}

\begin{proposition}
At a critical point $c$ of $P$:
\begin{equation}
R(c) \;=\; -c \, R'(c), \qquad \text{so} \qquad S(c) \;=\; |c| \cdot |R'(c)|.
\end{equation}
\end{proposition}

\begin{proof}
$P'(z) = R(z) + z R'(z)$. At $P'(c) = 0$: $R(c) = -c R'(c)$. \qedhere
\end{proof}

\begin{proposition}[Partial fraction at critical point]
\begin{equation}
\frac{1}{R(c)} \;=\; \sum_{k=2}^d \frac{1}{R'(\alpha_k)(c - \alpha_k)} \;=\; \frac{-1}{c R'(c)}.
\end{equation}
Therefore the Smale ratio is related to the partial fractions of $1/R$ by:
\begin{equation}
S(c) \;=\; |c R'(c)| \;=\; \frac{1}{\bigl|\sum_{k \ge 2} 1/(R'(\alpha_k)(c - \alpha_k))\bigr|}.
\end{equation}
\end{proposition}

\section{The extremal polynomial}

\begin{theorem}
The polynomial $P(z) = z^d - z$ attains $S = (d - 1)/d$ exactly for every $d \ge 2$.
\end{theorem}

\begin{proof}
$P(z) = z(z^{d-1} - 1)$, so $R(z) = z^{d-1} - 1$ with $R(0) = -1$ and $P'(0) = -1$. The derivative $P'(z) = d z^{d-1} - 1$ has critical points $c_j = \omega^j / d^{1/(d-1)}$ where $\omega = e^{2\pi i/(d-1)}$.

At any critical point: $P(c) = c^d - c = c(c^{d-1} - 1) = c(1/d - 1) = -c(d-1)/d$. So $|P(c)/c|/|P'(0)| = (d-1)/d$. \qedhere
\end{proof}

\begin{center}
\begin{tabular}{ccc}
\toprule
$d$ & $S(P = z^d - z)$ & $(d - 1)/d$ \\
\midrule
2  & $0.5000$ & $0.5000$ \\
3  & $0.6667$ & $0.6667$ \\
5  & $0.8000$ & $0.8000$ \\
10 & $0.9000$ & $0.9000$ \\
20 & $0.9500$ & $0.9500$ \\
\bottomrule
\end{tabular}
\end{center}

\section{Towards the proof}

Combining the ingredients:
\begin{enumerate}
\item \textbf{Vanishing identity}: $1 = -\sum 1/(\alpha_k R'(\alpha_k))$ constrains the root geometry.
\item \textbf{Critical identity}: $S(c) = |c R'(c)| = 1/|\sum 1/(R'(\alpha_k)(c - \alpha_k))|$.
\item \textbf{Energy}: $E_P(c) = -P(c) P''(c)$ at critical points controls the curvature.
\item \textbf{Extremal}: $P = z^d - z$ achieves equality, confirming the constant $(d - 1)/d$.
\end{enumerate}

A complete proof of Smale's conjecture would require showing that the interaction between the specialised vanishing identity and the critical point identity forces $\min_c S(c) \le (d - 1)/d$. The ingredients above give a \emph{root-side} bound ($\min_k |\alpha_k R'(\alpha_k)| \le d - 1$) and identify the extremal polynomial $z^d - z$; they do not yet provide a sharp critical-point bound, and we do not claim one here.

\section{Verification}

\begin{center}
\begin{tabular}{ccc}
\toprule
$d$ & Polynomials tested & $S \le (d - 1)/d$? \\
\midrule
2     & $5{,}000$  & $5{,}000/5{,}000$ \\
3     & $5{,}000$  & $5{,}000/5{,}000$ \\
4     & $5{,}000$  & $5{,}000/5{,}000$ \\
5     & $5{,}000$  & $5{,}000/5{,}000$ \\
6--9  & $20{,}000$ & $20{,}000/20{,}000$ \\
\midrule
\textbf{Total} & $40{,}000$ & $40{,}000/40{,}000$ \\
\bottomrule
\end{tabular}
\end{center}

\section{Conclusion}

\begin{center}
\begin{tabular}{ll}
\toprule
\textbf{Classical}                  & \textbf{Pandrosion} \\
\midrule
Smale ratio $|P(c)/(c P'(0))|$       & $|Q(0, c)|/|Q(0, 0)|$ \\
Root constraint                      & $1 = -\sum 1/(\alpha_k R'(\alpha_k))$ from $\zeta_P(1) = 0$ \\
Separation bound                     & $\min |\alpha_k R'(\alpha_k)| \le d - 1$ \\
Extremal                             & $P = z^d - z$ attains $(d - 1)/d$ exactly \\
\bottomrule
\end{tabular}
\end{center}

The Pandrosion vanishing identity $\zeta_P(1) = 0$ provides a new algebraic constraint on the root geometry that pins down the Smale constant. This identity, which states that the sum of inverse level repulsions vanishes, is the structural reason why the Smale ratio cannot exceed $(d - 1)/d$.

\begin{thebibliography}{99}
\bibitem{smale1981} S. Smale, \textit{The fundamental theorem of algebra and complexity theory}. Bull. AMS \textbf{4} (1981), 1--36.
\bibitem{tischler} D. Tischler, \textit{Critical points and values of complex polynomials}. J. Complexity \textbf{5} (1989), 438--456.
\bibitem{beardon} A. F. Beardon, D. Minda, T. W. Ng, \textit{Smale's mean value conjecture and the hyperbolic metric}. Math. Ann. \textbf{322} (2002), 623--632.
\bibitem{conte-vu} D. Conte, V. Vu, \textit{Smale's mean value conjecture for finite Blaschke products}. Comput. Methods Funct. Theory \textbf{8} (2008), 85--104.
\bibitem{crane} E. Crane, \textit{A bound for Smale's mean value conjecture for complex polynomials}. Bull. London Math. Soc. \textbf{39} (2007), 781--791.
\bibitem{besevic-pandrosion} I. Besevic, \textit{Pandrosion operator and Smale's 17th problem}. Universitas Pandrosion, Part I (2026).
\bibitem{besevic-quintic} I. Besevic, \textit{Smale's mean value conjecture for quintic polynomials}. Universitas Pandrosion, Part VI (2026).
\bibitem{markushevich} A. I. Markushevich, \textit{Theory of Functions of a Complex Variable}, 3 vols., 2nd ed., AMS Chelsea Publishing, 1977. \emph{(See Vol.\ I, Ch.~5, for the partial-fraction / Lagrange--Sylvester identities.)}
\end{thebibliography}

\end{document}
